\begin{figure}[htbp]
\centering
\begin{tikzpicture}[
    dot/.style={circle, fill=cNeutral!45, draw=none, inner sep=0pt, minimum size=1.7mm},
    hot/.style={circle, fill=cAlert, draw=none, inner sep=0pt, minimum size=2.3mm},
    sa/.style={-{Stealth[length=1.4mm,width=1.1mm]}, line width=0.5pt},
    jump/.style={line width=0.4pt, draw=faint, dash pattern=on 1.2pt off 1.2pt},
  ]
  \def\g{0.46}

  % ---------- (a) raster scan ----------
  \begin{scope}[shift={(0,0)}]
    \foreach \r in {0,...,3}{
      \foreach \c in {0,...,2}{
        \draw[sa, draw=cRGB] ({\c*\g+0.07},{-\r*\g}) -- ({(\c+1)*\g-0.07},{-\r*\g});
      }
      \ifnum\r<3 \draw[jump] ({3*\g},{-\r*\g}) -- (0,{-(\r+1)*\g}); \fi
    }
    \foreach \r in {0,...,3} \foreach \c in {0,...,3} \node[dot] at ({\c*\g},{-\r*\g}) {};
    \node[hot] at ({1*\g},{-1*\g}) {};
    \node[hot] at ({1*\g},{-2*\g}) {};
    \node[grp] at ({1.5*\g},0.62) {(a) raster};
    \node[lbl, text width=2.6cm] at ({1.5*\g},-2.25) {vertical neighbours\\ are four steps apart};
  \end{scope}

  % ---------- (b) bidirectional ----------
  \begin{scope}[shift={(3.35,0)}]
    \foreach \r in {0,...,3}{
      \foreach \c in {0,...,2}{
        \draw[sa, draw=cRGB] ({\c*\g+0.07},{-\r*\g+0.07}) -- ({(\c+1)*\g-0.07},{-\r*\g+0.07});
        \draw[sa, draw=cIR]  ({(\c+1)*\g-0.07},{-\r*\g-0.07}) -- ({\c*\g+0.07},{-\r*\g-0.07});
      }
    }
    \foreach \r in {0,...,3} \foreach \c in {0,...,3} \node[dot] at ({\c*\g},{-\r*\g}) {};
    \node[grp] at ({1.5*\g},0.62) {(b) bidirectional};
    \node[lbl, text width=2.6cm] at ({1.5*\g},-2.25) {forward and reverse,\\ as in Vim};
  \end{scope}

  % ---------- (c) four-way cross-scan ----------
  \begin{scope}[shift={(6.7,0)}]
    \def\h{0.3}
    % one mini grid per direction
    \begin{scope}[shift={(0,0)}]
      \foreach \a in {0,1,2} \foreach \b in {0,1}
        \draw[sa, draw=cRGB] ({\b*\h+0.05},{-\a*\h}) -- ({(\b+1)*\h-0.05},{-\a*\h});
      \foreach \a in {0,1,2} \foreach \b in {0,1,2} \node[dot, minimum size=1.3mm] at ({\b*\h},{-\a*\h}) {};
    \end{scope}
    \begin{scope}[shift={(1.15,0)}]
      \foreach \a in {0,1,2} \foreach \b in {0,1}
        \draw[sa, draw=cIR] ({(\b+1)*\h-0.05},{-\a*\h}) -- ({\b*\h+0.05},{-\a*\h});
      \foreach \a in {0,1,2} \foreach \b in {0,1,2} \node[dot, minimum size=1.3mm] at ({\b*\h},{-\a*\h}) {};
    \end{scope}
    \begin{scope}[shift={(0,-1.05)}]
      \foreach \a in {0,1} \foreach \b in {0,1,2}
        \draw[sa, draw=cFuse] ({\b*\h},{-\a*\h-0.05}) -- ({\b*\h},{-(\a+1)*\h+0.05});
      \foreach \a in {0,1,2} \foreach \b in {0,1,2} \node[dot, minimum size=1.3mm] at ({\b*\h},{-\a*\h}) {};
    \end{scope}
    \begin{scope}[shift={(1.15,-1.05)}]
      \foreach \a in {0,1} \foreach \b in {0,1,2}
        \draw[sa, draw=cTime] ({\b*\h},{-(\a+1)*\h+0.05}) -- ({\b*\h},{-\a*\h-0.05});
      \foreach \a in {0,1,2} \foreach \b in {0,1,2} \node[dot, minimum size=1.3mm] at ({\b*\h},{-\a*\h}) {};
    \end{scope}
    \node[grp] at (0.875,0.62) {(c) cross-scan};
    \node[lbl, text width=2.8cm] at (0.875,-2.25) {four directions, merged,\\ as in VMamba SS2D};
  \end{scope}

  % ---------- (d) local windows ----------
  \begin{scope}[shift={(10.25,0)}]
    \def\w{0.36}
    \foreach \wx/\wy in {0/0, 0.95/0, 0/-0.95, 0.95/-0.95}{
      \begin{scope}[shift={(\wx,\wy)}]
        \draw[sa, draw=cRGB] (0.06,0) -- ({\w-0.06},0);
        \draw[jump] (\w,0) -- (0,-\w);
        \draw[sa, draw=cRGB] (0.06,-\w) -- ({\w-0.06},-\w);
        \foreach \a in {0,1} \foreach \b in {0,1} \node[dot] at ({\b*\w},{-\a*\w}) {};
      \end{scope}
    }
    \node[grp] at (0.66,0.62) {(d) local windows};
    \node[lbl, text width=2.6cm] at (0.66,-2.25) {scan inside each window,\\ as in LocalMamba};
  \end{scope}
\end{tikzpicture}

\vspace{2mm}
\caption[Scan orders used to apply state-space models to images]{Scan orders
used to apply state-space models to images, shown on small grids of feature
cells. (a)~A raster scan places vertical neighbours (red) a full row apart
in sequence order. (b)~Vim adds a reverse pass. (c)~VMamba's SS2D scans rows
and columns in both directions and merges the four outputs. (d)~Local scans
keep neighbourhoods together. Work in 2026 replaces fixed orders with
adaptive or non-directional ones (Section~\ref{sec:vssm}).}
\label{fig:scans}
\end{figure}
